\emph{证明系统}提供了纯句法的方法来刻画语句间的语义后承与相容性。\emph{自然演绎}是其中一种证明系统。其中的一个\emph{推导}由一棵公式树构成。推导中最顶层的公式是\emph{假设}。为使推导正确，所有其他公式必须由若干\emph{推理规则}之一得到正确证成。这些规则是成对出现的：每个联结词和量词都有一个引入规则和一个消去规则。例如，如果一个公式~$!A$由 $\Elim{\lif}$ 规则证成，那么它前面的公式（即\emph{前提}）必须是 $!B \lif !A$ 和 $!B$（对某个~$!B$）。某些推理规则还允许将假设\emph{解除}。例如，如果使用 $\Intro{\lif}$ 从 $!B$ 推出 $!A \lif !B$，那么在推导出前提~$!B$的过程中作为假设出现的任意 $!A$ 都可以被解除，并且这些假设会被赋予一个标签，该标签也会记录在这个推理步骤处。

如果存在一个以~$!A$ 为结尾公式且所有假设都被解除的推导，我们就称 $!A$ 是一个\emph{定理}，并记作~$\Proves !A$。如果所有未解除的假设都在某个集合~$\Gamma$ 中，我们就称 $!A$ 是\emph{从~$\Gamma$可推导的}，并记作 $\Gamma \Proves !A$。如果 $\Gamma \Proves \lfalse$，我们称 $\Gamma$ 是\emph{不一致的}，否则是\emph{一致的}。这些概念是相互关联的，例如，$\Gamma \Proves !A$ 当且仅当 $\Gamma \cup \{\lnot !A\}$ 是不一致的。它们也与相应的语义概念相联系，例如，如果 $\Gamma \Proves !A$，那么 $\Gamma \Entails !A$。证明系统的这一性质——即凡是能从 $\Gamma$ 推导出的结论，都保证是~$\Gamma$的语义后承——被称为\emph{可靠性}。\emph{可靠性定理}通过对推导的长度进行归纳来证明，它表明：只要推理的前提被其未解除的假设所语义后承，每个单独的推理就保持其结论被未解除的假设所语义后承这一性质。